% document formatting
\documentclass[10pt]{article}
\usepackage[utf8]{inputenc}
\usepackage[left=1in,right=1in,top=1in,bottom=1in]{geometry}
\usepackage[T1]{fontenc}
\usepackage{xcolor}

% math symbols, etc.
\usepackage{amsmath, amsfonts, amssymb, amsthm}
\usepackage{dsfont}

% lists
\usepackage{enumerate}

% images
\usepackage{graphicx} % for images
\usepackage{tikz}
\usetikzlibrary{fit}

% code blocks
\usepackage{minted, listings} 

% verbatim greek
\usepackage{alphabeta}
\DeclareMathOperator*{\argmin}{arg\,min}
\DeclareMathOperator*{\argmax}{arg\,max}

\graphicspath{{./assets/images/}}

\title{02-750 Week 7 \\ \large{Automation of Scientific Research}}
 
\author{Aidan Jan}

\date{\today}

\begin{document}
\maketitle

\section*{IWAL Algorithm}
``Importance Weighted Active Learning''
\begin{itemize}
	\item Key Features
	\begin{itemize}
	    \item IWAL is an algorithm for learning \textbf{binary and multi-class classifiers}
	    \item It is a Type I algorithm (i.e., hypothesis elimination)
	    \item It can be used with any data access model
	    \item It can be used with \textbf{multiple loss functions} (Note: CAL, $A^2$, and DHM assume 0-1 loss)
	    \item It handles non-separable data
	    \item It is mellow and compensates for sampling bias using importance weighting
    \end{itemize}
	\item Guarantees
	\begin{itemize}
	    \item IWAL will converge to the optimal model
	    \item IWAL bounds the number of samples needed to reach the optimal model
    \end{itemize}   
\end{itemize}

\subsection*{IWAL Algorithm}
\begin{verbatim}
    Set S_0 = {}
    for t from 1, 2, ... until data stream runs out
        Receive x_t
        set p_t = rejection-threshold(x_t, history {x_i, y_i, p_i, Q_i : 1 <= i < t})
        flip a coin Q_t in {0, 1} with E[Q_t] = p_t.  
            If Q_t = 1, request y_t and set
                S_t = S_{t - 1} + {(x_t, y_t, p_min / p_t)}
            Else
                S_t = S_{t - 1}
        Let h_t = argmin_{h in H} sum_{x, y, c in S_i}(c * l(h(x), y))
\end{verbatim}
\begin{itemize}
	\item The arguments for IWAL include a subroutine that returns the probability of labeling $x_t$
	\begin{itemize}
	    \item For the purpose of computing importance weights, $p_t = P_{train}(x)$
	    \item $p_t > 0 \forall x_t$ and will be different for each $x_t$, in general
	    \item The paper presents two options for defining this subroutine; however, there are other options
    \end{itemize}
    \item The $p_{min}$ argument is required for technical reasons.
    \begin{itemize}
	    \item It is a constant $0 < p_{min} < 1$
    \end{itemize}
    \item The flip coin step will label $x_t$ with a probability $p_t$ randomly.  If the algorithm does label $x_t$, it adds it to set $S$ with importance weight $\frac{p_{min}}{p_t}$
    \item The last step learns an importance-weighted classifier, where $c = \frac{p_{min}}{p_t}$
    \begin{itemize}
	    \item The importance-weighted estimate of the loss after $T$ rounds is:
	    \[L_h(h) = \frac{1}{T} \sum_{t = 1}^T \frac{Q_t}{p_t} l(h(x_t), y_t)\]
        \item The authors show that $\mathbb{E}[L_T(h)] = L(h)$, refer to paper for details
    \end{itemize}
\end{itemize}

\subsection*{Rejection Threshold Subroutines}
\subsubsection*{Approach 1: Loss-weighting}
\begin{itemize}
	\item First, find the model that minimizes the importance-weighted loss
	\[L_{t - 1}^* = \min_{h \in H_{t - 1}} \frac{1}{t - 1} \sum_{i = 1}^{t - 1} \frac{Q_i}{p_i} l(h(x_i), y_i)\]
    \item Next, keep any model that isn't significantly worse than the best model
    \[H_t = \{h \in H_{t - 1} \::\: \frac{1}{t - 1} \sum_{i = 1}^{t - 1} \frac{Q_i}{p_i} l(h(x_i), y_i) \leq L_{t - 1}^* + \Delta_{t - 1}\}\]
    \begin{itemize}
	    \item $\Delta_{t - 1}$ is a generalization bound, like $A^2$ and DHM.
    \end{itemize}
    \item Finally, return $p_t = \max_{f, g \in H_t, y \in Y} l(f(x), y) - l(g(x), y)$
\end{itemize}
The rejection probability for $x$ is the largest difference in losses between any pair of models in the version space, considering each potential label for $x$.\\\\
A large difference reflects a significant disagreement between the models, and thus an opportunity to reduce the size of the version space

\subsubsection*{Approach 2: Bootstrapping}
This subroutine is easier to implement than the previous one.
\begin{itemize}
	\item Parameters: $b$ (length of initial boostrap sample), $k$ (number of predictors used in voting), $p_{min}$ (lower bound on rejection threshold)
	\item First, if $t \leq b$, set $p_t = 1$
	\item If $t = b$, train predictors $h_1, \dots, h_k$ on the initial sample $(x_1, y_1, c_t), \dots, (x_b, y_b, c_t)$.  Denote the set by $H$.  (This is constructing a committee.)
	\item If $t > b$, set $p_t = p_{min} + (1 - p_{min})[\max_{h_i, h_j \in H, y \in Y} L(h_i(x), y) - L(h_j(x), y)]$
	\item Return $p_t$
\end{itemize}

\subsubsection*{Approach 3: Critical Importance Weight}
This subroutine is even easier to implement than the previous one.
\begin{itemize}
	\item Let $h = Learn(S)$ be the importance-weighted model trained on $S$.
	\item Let $h' = Learn(S \cup \{x, \neg h(x), w_\Delta\})$ be a model trained with $x$ assigned to some label other than the one predicted by $h$, and using importance weight $w_{\Delta}$
	\item Search for $w_\Delta^*$, which is the smallest $w_\Delta$ such that $h(x) \neq h'(x)$
	\item Set $p_t = \frac{1}{w_\Delta^*}$
\end{itemize}
You can think of $w_\Delta^*$ as being the amount of 'energy' needed to create a model $h'$ that would assign $x$ to a different label than model $h$.  If $w_\Delta^*$ is small, then it must be easy to find a model that disagrees with $h$.  If $w_\Delta^*$ is large, then it must be hard to find a model that disagrees with $h$.

\subsection*{IWAL Guarantees}
\begin{itemize}
	\item Convergence (Assuming that $p_t > 0, \forall x \in \mathcal{X}$)
	\begin{itemize}
	    \item IWAL will return a near-optimal model, with probability at least $1 - \delta$
	    \item Notice that we avoid the consistency issue from earlier, even though we have a biased set of labels because $p_t > 0$
    \end{itemize}
	\item Sample Complexity
	\begin{itemize}
	    \item Like $A^2$ and DHM, the label complexity depends on the amount of noise
	    \item The key result is that IWAL is at worst a constant factor than standard supervised learning (i.e., passive learning) in terms of sample complexity
	    \item The empirical results suggest that IWAL is often much better.
    \end{itemize}
\end{itemize}

\subsection*{Summary}
$A^2$ and DHM are essentially versions of CAL for non-separable data
\begin{itemize}
	\item $A^2$ significance is mainly that it was the first active learning algorithm to be able to handle non-separable data and still provide guarantees with respect to convergence and label complexity
	\begin{itemize}
	    \item Unfortunately, it can't be implemented exactly (for most hypothesis spaces)
	    \item Running it using sampling is computationally expensive
    \end{itemize}
    \item DHM also handles non-separable data and comes with guarantees
    \begin{itemize}
	    \item Practical and relatively easy to implement
    \end{itemize}
\end{itemize}
IWAL also handles non-separable data and comes with guarantees
\begin{itemize}
	\item Like $A^2$ and DHM, the label complexity depends on the amount of noise
	\item IWAL is at worst a constant factor than standard supervised learning (i.e., passive learning) is terms of sample complexity
	\item The empirical results suggest that IWAL is often much better
\end{itemize}

\section*{The ZLG Algorithm}
Key features:
\begin{itemize}
	\item ZLG is an algorithm for learning \textbf{binary classifiers}
	\item It is a Type II algorithm (i.e., cluster exploitation)
	\item It can be used with \textbf{pool-based selective sampling}
	\item It is \textbf{aggressive}
	\begin{itemize}
	    \item It approximately minimizes expected risk
    \end{itemize}
    \item Guarantees: None
\end{itemize}

\subsection*{Assumptions}
\begin{itemize}
	\item \textbf{Key assumption:} if two unlabeled instances, $x_i$ and $x_j$, are similar, then they should have the same label
	\item There are many ways we might measure similarity.  The ZLG paper uses a \textbf{radial basis function}.
	\begin{itemize}
	    \item Let $x \in \mathbb{R}^m$.  The similarity is measured as follows:
	    \[w_{ij} = \exp \left(-\frac{1}{\sigma^2} \sum_{d = 1}^m (x_i(d) - x_j(d))^2\right)\]
        $w_{ij}$ will be large if $x_i$ and $x_j$ are similar, otherwise, $w_{ij}$ will be small
    \end{itemize}
\end{itemize}

\subsection*{Algorithm}
ZLG builds a graph over all instances in the pool using weights
\begin{itemize}
	\item Let $G = (V, E, W)$ be a weighted graph where:
	\begin{itemize}
	    \item There is one node for each of the $n$ instances in the pool
	    \item Nodes $i$ and $j$ are connected by and edge if $w_{ij} \geq t$
	    \item Here, $t$ is a user-defined threshold
	    \item $W$ is a $n \times n$ weight matrix such that
	    \[W(i, j) = \begin{cases} w_{ij} & \text{if } w_{ij} \geq t \\ 0 & \text{otherwise} \end{cases}\]
    \end{itemize}
\end{itemize}
ZLG then considers the labels of the nodes in the graph.
\begin{center} 
	\includegraphics*[width=0.9\textwidth]{W7_1.png} 
\end{center}
\begin{itemize}
	\item Labels are binary (0 or 1)
	\begin{itemize}
	    \item This implies there are $2^n$ possible joint labelings on a graph with $n$ nodes
	    \item But, under our assumption that if unlabeled instances $x_i$ and $x_j$ are similar, they should have the same label, \textbf{some labelings are more likely than others}
    \end{itemize}
\end{itemize}
We can formlaize our assumption by defining a model for the joint probability distribution over the labels, $P(Y_1, \dots, Y_n)$
\begin{itemize}
	\item 'Good' labelings will have high probability under $P(Y_1, \dots, Y_n)$
	\item The most obvious choice for specifying $P(Y_1, \dots, Y_n)$ is as a multivariate \textbf{binomial} (i.e., a set of coupled binomial variables), but the authors use a \textbf{multivariate Gaussian distribution} instead
	\begin{itemize}
	    \item This is done for mathematical/computational convenience and tractablility
	    \item This simplification implies that each $Y_i$ is now continuous and that its marginal distribution is Gaussian
	    \item If we are asked to make the prediction for some particular $Y_i$, we first infer its value (under the multivariate Gaussian) and then round to 0 or 1
    \end{itemize}
\end{itemize}

\subsubsection*{Aside: Multivariate Gaussian Distributions}
\begin{itemize}
	\item Let $Y = \{Y_1, \cdots, Y_n\}$ be a set of $n$ Gaussian random variables
	\item Let $\mu = \{\mu_1, \cdots, \mu_n\}$ be a mean vector over $Y$
	\item Let $\Sigma$ be the covariance matrix over $Y$
	\item Let $\Sigma^{-1}$ be the precision matrix over $Y$
	\item The probability density of a vector $y \in \mathbb{R}^n \sim \mathcal{N}(\mu, \Sigma)$ is:
	\[P(y) = \frac{1}{\sqrt{(2\pi)^n |\Sigma|}} \exp \left(-\frac{1}{2} (y - \mu)^T \Sigma^{-1} (y - \mu)\right)\]
    where $|\Sigma|$ denotes the \textbf{determinant} of $\Sigma$
\end{itemize}

\subsection*{Overview of the ZLG Algorithm}
\begin{itemize}
	\item Let $Y_u \subseteq Y = \{Y_1, \cdots, Y_n\}$ be the set of \textbf{unknown labels}
	\item Let $Y_k \subseteq Y = \{Y_1, \cdots, Y_n\}$ be the set of \textbf{known labels}.  Naturally, $Y_u \cap Y_k = \emptyset$
	\item The ZLG algorithm uses the graph and any known labels to construct the probability distribution, $P(Y_u | Y_k)$
	\item It then solves for $\mu_u$
	\begin{itemize}
	    \item The elements of $\mu_u$ can be rounded to 0 or 1 to make predictions
	    \item Alternatively, use the elements of $\mu_u$ to select the next instance to query
    \end{itemize}
\end{itemize}
ZLG defines $P(Y_u | Y_k)$ as a \textbf{Boltzmann distribution}
\begin{itemize}
	\item That is, the probability of any vector $y$ is inversely proportional to the user-defined \textbf{energy function}
	\item The energy used by the ZLG algorithm is:
	\[E(y) = \frac{1}{2} \sum_{i, j} w_{ij} (y_i - y_j)^2\]
    \item This is a \textbf{harmonic potential}.  It penalizes in proportion to the product of the edge weight (if any) and the square of the differences in labels
    \item This energy can be rewritten in matrix form:
    \begin{itemize}
	    \item Recall that $W$ is the $n \times n$ weight matrix
	    \item Let $D$ be a diagonal matrix such that $d_i = \Sigma_j w_{ij} = \Sigma_j W(i, j)$.  Thus, $d_i$ is the weighted degree of node $Y_i$
	    \item Let $\Delta = D - W$ be the Laplacian Matrix
	    \item The energy in matrix form is $E(y) = y^T \Delta y$
    \end{itemize}
\end{itemize}
It can be shown that the following is a multivariate Gaussian
\[P(y) = \frac{1}{Z_\beta}\exp (-\beta E(y))\] 
where $\beta$ is the `inverse temperature' parameter, and
\[Z_\beta = \int \exp(-\beta E(y)) \:\dd y\]
is the partition function (i.e., normalization constant)
\begin{itemize}
	\item Notice that:
	\begin{itemize}
	    \item The lowest energy will have the highest probability (Boltzmann dist.)
	    \item In the absence of \textit{any} known labels, there are two minimum-energy solutions, one where everything is labeled 0, and one where everything is labeled 1.
	    \item The formal name for their model is a \textbf{Gaussian Random Field}, which is a \textbf{probabilistic graphical model} (PGM) for encoding multivariate Gaussian distributions
    \end{itemize}
\end{itemize}
Suppose that we do have some labels
\begin{itemize}
	\item Conceptually, we could imagine a `wave' of 0's (resp. 1's) diffusing throughout the weighted graph emanating from the instances labeled 0 (resp. 1's)
	\item When these waves meet, they will combine to create some intermediate label between 0 and 1, depending on how far away they are from known labels.
\end{itemize}
\begin{center} 
	\includegraphics*[width=\textwidth]{W7_2.png} 
\end{center}

\subsection*{Implementation of ZLG}
It can be shown that this diffusion process can be implemented efficiently in a Gaussian Random Field in order to infer labels.
\begin{itemize}
	\item Reorganize the Laplacing Matrix as follows:
	\[\Delta = \begin{bmatrix} \Delta_{ll} & \Delta_{lu} \\ \Delta_{ul} & \Delta_{uu} \end{bmatrix}\]
    where $\Delta_{ll}, \Delta_{ul}, \Delta_{lu}, \Delta_{uu}$ are the submatrices corresponding to the groups of labeled (l) and unlabeled instances (u).
    \item The minimum-energy solution for the unlabeled instances can be computed as $f_u = -\Delta_{uu}^{-1} \Delta_{ul} f_l$, where $f_l$ is the vector of known labels
    \item Also, the estimated distribution over the unlabeled instances is:
    \[P(y_u | f_l) = \mathcal{N} (f_u \Delta_{uu}^{-1})\]
    \item Thus, we can compute the most likely labels using matrix operations!
\end{itemize}

\subsection*{Query Selection}
\begin{itemize}
	\item Notice that the predicted labels \textit{could} be used to perform uncertainty sampling, but we already know the problems with that approach
	\item Instead, the authors perform expected risk minimization
	\begin{itemize}
	    \item Let $f = [f_l, f_u]$ be the minimum energy solution
	    \item Let $\hat{R}(f) = \sum_{i = 1}^n \mathbb{I}[round (f_i) \neq 0](1 - f_i) + \mathbb{I}[round(f_i) \neq 1] f_i = \sum_{i = 1}^n \min(f_i, 1 - f_i)$ be the \textbf{estimated risk}, where $round(f_i)$ is the result of rounding to 0 or 1.
    \end{itemize}
    \item Recall that we round each $f_i$ to 0 or 1.  Thus, if $f_i > 0.5$, we predict node $i$ is 1.  $f_i$ is, therefore, our confidence that node $i$'s label is 1.  The estimated risk (i.e., generalization error) is the sum of the probabilities that the rounding selects the wrong label.
    \item If we request the label for node $k$ and we learn that its label, $y_k$, the updated estimated risk is:
    \[\hat{R}(f) = \sum_{i = 1}^n \min \left(f_i^{+(x_k, y_k)}, 1 - f_i^{+(x_k, y_k)}\right)\]
    \item Since we do not know which label the oracle will return, we compute the expected estimated risk if we were to request the label for node $k$
    \[\hat{R}(f^{+x_k}) = (1 - f_k) \hat{R} (f^{+(x_k, 0)}) + f_k \hat{R}(f^{+(x_k, 1)})\]
    \item The algorithm queries the node satisfying: $k = \argmin_{k'} \hat{R} (f^{+x_{k'}})$
    \item Note, computing $\argmin_{k'} \hat{R}(f^{+x_{k'}})$ can be done relatively efficiently due to properties of Gaussians (Refer to the paper for details.)
\end{itemize}

\subsection*{Pseudocode}
\begin{verbatim}
Build a neighborhood graph over data in the pool
Convert the neighborhood graph into a Gaussian Random Field 
Iteratively
    Infer labels using the model
    Greedily select the point expected to minimize future risk
    Call oracle to get true label for selected point
    Update the model
\end{verbatim}

\subsection*{Limitations of the ZLG algorithm}
\begin{itemize}
	\item The algorithm requires $O(n^3)$ work per iteration because it performs matrix inversion, which could be a problem for large $n$
	\item Sampling bias:
	\begin{center} 
	    \includegraphics*[width=0.8\textwidth]{W7_3.png} 
    \end{center}
    \item Suppose we have these observations.  The inference procedure effectively implements a diffusion process, and the most uncertain points will lie at or near the mid-point.  The model is over-confident and will continue to suffer from sample bias until a point from region $B$ is selected.
\end{itemize}





\end{document}